\documentclass{situationist-preprint}

\title{Happiness as a By-Product}
\msid{EXI-07-a}
\msfamily{EXI-07}
\msdiscipline{Meaning, Death and the Absurd}
\mstype{Principal manuscript}
\mskeywords{happiness; paradox of hedonism; by-products; Sidgwick; Mill; Butler; final ends}

\begin{document}
\maketitle

\begin{abstract}
Is happiness the goal of life? A long tradition answers yes: whatever else we pursue, we pursue for the sake of being happy. We argue that the answer is no. Happiness is a by-product of other pursuits, and aiming at it directly is self-defeating. The argument develops the paradox of hedonism, stated by Sidgwick, anticipated by Butler and recorded by Mill from his own crisis: the pleasures of pursuit depend on caring about objects other than pleasure. We give the paradox a precise form. Happiness, on any plausible account, consists to a large extent in engagement with and success in activities whose objects are valued for their own sake. An agent whose only final end is her own happiness values nothing else for its own sake, and so lacks the very attitudes on which most of her happiness depends. The policy of making happiness one's sole final end therefore undermines its own aim. We distinguish this thesis from the weaker claim that monitoring one's happiness is distracting, answer the reply that a sophisticated hedonist can adopt other ends instrumentally, and conclude that happiness is a good of life but not its goal.
\end{abstract}

\section{The problem}

Aristotle observed that everyone agrees on the name of the highest good, which is happiness, and that we choose it always for itself and never for the sake of anything else, while we choose honor, pleasure and understanding both for themselves and for the sake of happiness (\emph{Nicomachean Ethics} I.4, I.7). Bentham and the classical utilitarians made a narrower claim about the same structure: pleasure and the absence of pain are the only things desirable as ends, and everything else is desirable as a means to them. On both views happiness is \emph{the goal of life} in a strong sense: it is the end that all other ends serve, and a life is well directed to the extent that it is directed at it.

Against this view stands an observation as old as the view itself. \citet[Sermon XI]{Butler1726} argued that self-love, the general desire for one's own happiness, cannot itself be the source of our enjoyments, since enjoyment consists in the gratification of particular passions, for food, for esteem, for the good of others, whose objects are things other than our happiness. A being with self-love alone would have nothing to enjoy. \citet[I.4 and II.3]{Sidgwick1907} took up Butler's point and named its consequence the fundamental paradox of egoistic hedonism: the impulse toward pleasure, if too predominant, defeats its own aim. \citet[ch.~V]{Mill1873} reports that he recovered from the depression of his twentieth year by a change in his theory of life: those only are happy, he came to think, who have their minds fixed on some object other than their own happiness, such as the happiness of others or some art or pursuit followed as an ideal end; to ask oneself whether one is happy is to cease to be so. \citet[ch.~II]{Elster1983} generalized the structure: some states, such as spontaneity, sleep and the impression one makes on others, are \emph{essentially by-products}, which cannot be brought about intelligently and intentionally.

Recent empirical work suggests that the phenomenon is real. \citet{MaussEtAl2011} found that people who value happiness highly report less happiness and more depressive symptoms, and that inducing participants to value happiness reduced their enjoyment of a pleasant film. These findings are compatible with several explanations, and we do not rest the argument on them. Our task is to state the paradox so that it supports the conclusion of this family: happiness is not the goal of life, because it is a by-product of other pursuits, and aiming at it directly is self-defeating.

\section{Framework}

We use \enquote{happiness} broadly, for a life's going well for the one who lives it as she experiences it: the predominance of enjoyment, satisfaction and positive emotional condition over their opposites. The argument does not depend on a choice between hedonistic, life-satisfaction and emotional-state accounts \citep{Haybron2008}; it depends only on what makes people happy, on any of these accounts.

\begin{definition}[Final end]\label{d:final}
An agent has $X$ as a \emph{final end} if she desires $X$ for its own sake. $X$ is her \emph{sole} final end if she desires everything else she desires only as a means to $X$.
\end{definition}

\begin{definition}[Direct aim]\label{d:direct}
An agent \emph{aims directly} at happiness if her own happiness is her sole final end, so that her practical reasoning selects activities, relationships and commitments by their expected contribution to it.
\end{definition}

\begin{definition}[By-product]\label{d:byprod}
A good $G$ is a \emph{by-product} of a class of pursuits if $G$ is realized to a large extent through those pursuits, and the pursuits contribute to $G$ only when the agent values their objects for their own sake.
\end{definition}

\begin{principle}[Engagement]\label{p:eng}
A large part of human happiness consists in, or depends causally on, wholehearted engagement in activities and relationships: absorption in work, play and inquiry, the enjoyment of others' company and good, the satisfaction of achieving what one set out to do.
\end{principle}

\begin{principle}[Object-directedness]\label{p:obj}
The forms of engagement in Principle~\ref{p:eng} require that the agent care about their objects for their own sake: that the outcome of the game matter to the player, the truth to the inquirer, the friend's good to the friend, independently of the satisfaction that caring yields.
\end{principle}

Principle~\ref{p:eng} is the common ground of the psychology of well-being and of every theory of happiness that does not reduce it to bodily pleasure. Principle~\ref{p:obj} is Butler's point. It is a claim about the structure of the attitudes, not about their phenomenology. The satisfaction of having solved a problem is satisfaction \emph{that the problem is solved}; it presupposes that its being solved matters, and an agent to whom it matters only as a source of satisfaction has nothing for the satisfaction to be about. The pleasure of a friend's company is in part pleasure in the friend, and it is not available to someone who treats the friend as an instrument for producing pleasure in himself, both because his attitude is not friendship and because friends notice.

\section{The argument}

\begin{lemma}\label{l:byprod}
Happiness is a by-product of engagement in activities and relationships valued for their own sake.
\end{lemma}

\begin{proof}
By Principle~\ref{p:eng}, happiness is realized to a large extent through engagement. By Principle~\ref{p:obj}, engagement contributes to happiness only when the agent values its objects for their own sake. These are the two clauses of Definition~\ref{d:byprod}.
\end{proof}

\begin{lemma}\label{l:exclude}
An agent who aims directly at happiness does not value the objects of her activities for their own sake.
\end{lemma}

\begin{proof}
By Definition~\ref{d:direct}, her happiness is her sole final end, so by Definition~\ref{d:final} she desires everything else only as a means to it. To desire the outcome of a game, the truth of a theorem or the good of a friend only as a means to one's happiness is not to value it for its own sake.
\end{proof}

\begin{theorem}\label{t:main}
Aiming directly at happiness is self-defeating: an agent who does so forgoes the large part of happiness that depends on engagement, and attains less happiness than she would if she had final ends other than her happiness.
\end{theorem}

\begin{proof}
By Lemma~\ref{l:exclude}, the agent who aims directly at happiness does not value the objects of her activities for their own sake. By Lemma~\ref{l:byprod}, the part of happiness that is realized through engagement is available only to agents who do. She therefore forgoes that part, which by Principle~\ref{p:eng} is large. An agent who values other things for their own sake has access to it, and has access to whatever remains of happiness besides. The policy of direct aim thus yields less of its own target than an available alternative, which is what it is for a policy to be self-defeating.
\end{proof}

\begin{corollary}\label{c:goal}
Happiness is not the goal of life, in the sense of the end that all other ends of a well-directed life serve.
\end{corollary}

\begin{proof}
If happiness were the goal of life in that sense, a well-directed life would have happiness as its sole final end, and so would aim directly at it. By Theorem~\ref{t:main}, such a life defeats the aim by which it is directed. A life directed at its goal in a way that defeats that goal is not well directed. So a well-directed life does not have happiness as its sole final end.
\end{proof}

The theorem concerns final ends, not attention. It is sometimes taken to say only that thinking about one's happiness spoils it, as Mill's maxim about asking oneself whether one is happy suggests. That weaker thesis is true but contingent: an agent could perhaps aim at happiness while keeping the aim out of view. Theorem~\ref{t:main} does not depend on attention. An agent who never once thinks about her happiness, but for whom everything else is a means to it, still lacks the attitudes that Principle~\ref{p:obj} requires.

Corollary~\ref{c:goal} fits the plural account of the good defended in ETH-01, on which well-being is one of seven irreducible components of the good alongside knowledge, achievement, relationships, fairness, beauty and autonomy, and with ETH-05's argument that well-being itself includes contact with reality. If knowledge, achievement and relationships are good in themselves, valuing them for their own sake is not a psychological device for being happy but a correct response to their value, and the happiness that follows is the by-product of getting their value right. It fits also the account of meaning in EXI-01 and EXI-17, on which meaning comes from active engagement with what is worthwhile.

\section{Objections and replies}

\begin{objection}[The sophisticated hedonist]\label{o:soph}
A clever egoist will see the paradox and adopt other final ends because doing so makes her happier. Her ultimate end remains her happiness, and her policy is not self-defeating. This is the structure of the sophisticated consequentialist of \citet{Railton1984}.
\end{objection}

\begin{reply}
If she succeeds, she comes to desire the objects of her pursuits for their own sake, and then her happiness is no longer her sole final end (Definition~\ref{d:final}). The sophisticated hedonist does not aim directly at happiness in the sense of Definition~\ref{d:direct}; she has replaced that aim, for reasons that refer to happiness, with others. This confirms the thesis: the path to happiness runs through ceasing to treat it as the only end. If instead she keeps her happiness as the sole end and merely acts as if she valued other things, she has not acquired the attitudes on which engagement depends, and Theorem~\ref{t:main} applies to her. Railton's agent can stably value other things for their own sake only because his commitment to the further end is not a standing filter on every particular concern; once it is not, happiness is not functioning as the goal that all other ends serve.
\end{reply}

\begin{objection}[Aristotle]\label{o:arist}
For Aristotle, happiness is activity of the soul in accordance with virtue, and the virtuous person chooses virtuous acts for their own sake. Aiming at happiness in this sense is aiming at virtuous activity, and it is not self-defeating.
\end{objection}

\begin{reply}
Aristotle's \emph{eudaimonia} is not happiness in our sense but the good life as a whole, which includes the goods that a plural account counts as components of the good. Aiming at it is aiming at excellent activity in work, friendship and thought for its own sake, which is not aiming directly at happiness as Definition~\ref{d:direct} defines it. The Aristotelian claim that the good life is the end of action is compatible with our thesis, and indeed makes the same point: what is chosen for its own sake within such a life is not the agent's feeling good, but the activities, which the virtuous person values in themselves.
\end{reply}

\begin{objection}[Some happiness is aimed at]\label{o:aim}
People do aim at happiness successfully: they take holidays, eat well, avoid stress, seek therapy. These are direct pursuits of happiness, and they work.
\end{objection}

\begin{reply}
These are local pursuits of particular goods of experience, embedded in a life whose final ends lie elsewhere, and the thesis does not deny that they succeed. Theorem~\ref{t:main} concerns the policy of making one's happiness the \emph{sole} final end. The holiday helps because one returns to commitments one cares about; therapy for depression aims at restoring the capacity to care about things, which, on Mill's report, is what his recovery consisted in.
\end{reply}

\begin{objection}[Engagement is only causally required]\label{o:causal}
Perhaps engagement is merely a reliable cause of happiness. A technology that produced the feelings of engaged satisfaction without any engagement would then give the direct aim its target.
\end{objection}

\begin{reply}
On a pure mental-state account of happiness, the objection identifies a logically possible route to happiness that does not pass through other ends. Two points answer it. First, no such route is available to human beings, and the question whether happiness is the goal of life concerns the lives human beings can lead. Second, the experience machine shows that a life of produced satisfaction is worse for the one who lives it than a life of real engagement with the same experiences (ETH-05). Even if such a route made the direct aim succeed on its own terms, it would show that a life whose goal is happiness so understood is not the best life for its subject, which is a further reason to deny that happiness is the goal of life.
\end{reply}

\section{What the result does not establish}

The argument does not show that happiness is unimportant, or that one should be indifferent to it in deciding how to live. Happiness is a good, and it is a relevant consideration in choosing a career, a place to live or a partner. Nor does it show that happiness is not good for its own sake: a by-product can be intrinsically valuable. What it shows is that happiness cannot coherently be the end that organizes all the others. It also does not establish which other ends a person should have; that is the work of the account of the good (ETH-01) and of meaning (EXI-01, EXI-17).

\section{Conclusion}

Happiness is not the goal of life. Much of it consists in, or depends on, engagement with activities and people valued for their own sake, and an agent whose only final end is her own happiness values nothing else for its own sake. She therefore lacks the attitudes on which most of her happiness depends, and her aim defeats itself. Happiness is a by-product of caring about other things, as Butler argued, Sidgwick named and Mill discovered.

\begin{thebibliography}{99}\small
\bibitem[Butler(1726)]{Butler1726} Butler, J. (1726). \emph{Fifteen Sermons Preached at the Rolls Chapel}. London.
\bibitem[Elster(1983)]{Elster1983} Elster, J. (1983). \emph{Sour Grapes: Studies in the Subversion of Rationality}. Cambridge University Press.
\bibitem[Haybron(2008)]{Haybron2008} Haybron, D.~M. (2008). \emph{The Pursuit of Unhappiness: The Elusive Psychology of Well-Being}. Oxford University Press.
\bibitem[Mauss et~al.(2011)]{MaussEtAl2011} Mauss, I.~B., Tamir, M., Anderson, C.~L. and Savino, N.~S. (2011). Can seeking happiness make people unhappy? Paradoxical effects of valuing happiness. \emph{Emotion}, 11(4), 807--815.
\bibitem[Mill(1873)]{Mill1873} Mill, J.~S. (1873). \emph{Autobiography}. Longmans, Green, Reader, and Dyer.
\bibitem[Railton(1984)]{Railton1984} Railton, P. (1984). Alienation, consequentialism, and the demands of morality. \emph{Philosophy \& Public Affairs}, 13(2), 134--171.
\bibitem[Sidgwick(1907)]{Sidgwick1907} Sidgwick, H. (1907). \emph{The Methods of Ethics}, 7th edition. Macmillan.
\end{thebibliography}

\end{document}
